\section{\texorpdfstring{熵与散度}{熵与散度}}\label{app2-entropy-12-01532}
熵 \(H(\overline{\mathbf{p}})\) 在均匀分布处取得最大值。因此，给定一个散度函数，可以定义熵 \(H(\overline{\mathbf{p}})\) 为

\setcounter{equation}{109}
\begin{equation}
\label{FD107-entropy-12-01532}
\fitmath{H{(\overline{\mathbf{p}})}{= D(}\overline{\mathbf{p}}{\,||\, 1)} + c}
\end{equation}

其中 c 为某个常数，\(\overline{p}\) 为归一化概率测度。


在许多情形中，\(\Psi{(\overline{\mathbf{p}})}{= - D(}\overline{\mathbf{p}}{\,||\, 1)}\) 是 \(\overline{p}\) 的凸函数。例如，由 KL 散度可得

\setcounter{equation}{110}
\begin{equation}
\label{FD108-entropy-12-01532}
\fitmath{H_{KL}{(\overline{\mathbf{p}})} = - \int\overline{p}{(\mathbf{x})}\log\overline{p}{(\mathbf{x})}d\mu{(\mathbf{x})}}
\end{equation}

忽略一个常数后，这就是 Shannon 熵。


定义概率的幂函数为

\setcounter{equation}{111}
\begin{equation}
\label{FD109-entropy-12-01532}
\fitmath{h_{\alpha}{(\overline{\mathbf{p}})} = \int{\overline{p}}^{\alpha}{(\mathbf{x})}d\mu{(\mathbf{x})}}
\end{equation}

其中 \(0 < \alpha < 1\)。


由 \(\alpha\) 散度 \(D_{A}^{(\alpha)}\) 和 \(\beta\) 散度 \(D_{B}^{(\beta)}\)，得到相应的熵

\setcounter{equation}{112}
\begin{equation}
\label{FD110-entropy-12-01532}
\fitmath{H{(\overline{\mathbf{p}})} = c\, h_{\alpha}{(\overline{\mathbf{p}})} + d}
\end{equation}

其中 c 和 d 为常数，且 \(\alpha = \beta\)。


类似地，由 \(\gamma\) 散度 \(D_{G}^{(\gamma)}\) 得到

\setcounter{equation}{113}
\begin{equation}
\label{FD111-entropy-12-01532}
\fitmath{H{(\overline{\mathbf{p}})} = c\,\log{(h_{\alpha}{(\overline{\mathbf{p}})})} + d}
\end{equation}

其中 \(\alpha = \gamma\)。


不同的散度函数可以生成相似的熵函数族。反过来，从熵函数生成散度也有多种方法。一种典型方法，是利用凸生成函数 \(\Phi{(\overline{p})} = - H{(\overline{p})}\) 导出 Bregman 散度：

\setcounter{equation}{114}
\begin{equation}
\label{FD112-entropy-12-01532}
\fitmath{D_{\Phi}(\overline{\mathbf{p}}{\,||\,}\overline{\mathbf{q}}{) = \int}\left( {\Phi{(\overline{p})} - \Phi{(\overline{q})} - \frac{\delta\Phi(\overline{q})}{\delta\overline{q}}{(\overline{p} - \overline{q})}} \right)\, d\mu{(\mathbf{x})}}
\end{equation}
